\section{De clientes extremos a una plataforma general: Inductive Productization}

En el software empresarial ocurre con frecuencia una situación contraria a la intuición: que un solo cliente plantee cierto problema no significa que el problema carezca de importancia; ese cliente puede llegar a cierta restricción cinco años antes que el mercado. El ponente describe la enterprise demand como una power law, no como necesidades que se distribuyen alrededor de un promedio. Los equipos forward-deployed primero resuelven los casos extremos dentro del workflow real, después identifican las invariant a partir de los casos repetidos y, por último, construyen la interfaz del producto.

\subsection{No es desarrollo a la medida ni planeación de producto alejada del terreno}

El desarrollo puramente a la medida plasma cada diferencia entre clientes en una rama distinta, por lo que no permite acumular nada; la planeación de producto alejada del terreno, en cambio, tiende a promediar las restricciones extremas, que son las más valiosas. La Inductive productization se ubica entre ambas: la primera vez se registra el contexto completo, la segunda se buscan similitudes estructurales y solo a la tercera se empieza a formar una abstraction estable. El ponente atribuye la respuesta rápida en Operation Warp Speed a que antes ya se habían resuelto, en la optimización de la producción de petróleo y gas, problemas de coordinación de recursos con una estructura similar.

\begin{figure}[H]
\centering
\includegraphics[width=0.94\textwidth]{images/07-inductive-productization.png}
\caption{De un único cliente del futuro a un producto general: repetir, separar los detalles locales y codificar las invariant.}
\figsource{Redibujada a partir de los subtítulos 19:00--21:30, `images/07-inductive-productization.png`.}
\end{figure}

\begin{knowledgebox}{Lectura de la figura: qué cuenta como ``structurally identical''}
Dos industrias pueden usar sustantivos completamente distintos y, aun así, tener el mismo grafo de restricciones. Por ejemplo, si ambas tienen recursos escasos, prioridades de demanda, oferta dinámica, permisos de revisión y retroalimentación de resultados, pueden compartir primitive de optimization y de workflow. La similitud no puede basarse solo en analogías narrativas: hay que establecer de forma explícita la correspondencia entre entity, constraint, decision y outcome.
\end{knowledgebox}

\begin{warningbox}{La experiencia forward-deployed también puede sobreajustarse}
Los equipos en terreno tienden a tomar el proceso del cliente que más se hace oír como si fuera un hecho del mercado. Para evitar el sobreajuste hay que registrar qué lógica corresponde a regulaciones, a costumbres de la organización o a workaround temporales, y cuál constituye una verdadera invariant común a varios clientes. Después de la conversión en producto, también hay que verificar si los nuevos clientes pueden adoptarlo sin replicar el personal del equipo original.
\end{warningbox}

\subsection{Resumen de la sección}

Los entornos de tareas extremos pueden revelar de antemano problemas de infraestructura futuros, pero solo se convierten en una plataforma mediante la repetición, la comparación y la abstracción. Apollo/Rubix es en sí mismo el resultado de este camino: primero resolvió la entrega y la operación de la fleet interna y después se externalizó gradualmente como una capacidad general. El siguiente capítulo aplica la misma idea a Project Maven, que pasa de la detección visual a la cadena de decisión completa.
